\documentclass[11pt]{article}
\usepackage[margin=1in]{geometry}
\usepackage{amsmath,amssymb,amsthm}
\newtheorem{proposition}{Proposition}
\newtheorem{remark}{Remark}
\title{Band-limited Vorticity Direction: Corrected Estimates and a Sharpness Obstruction}
\author{Working correction for Jonathan Robert Simons\\Prime Field Technologies LLC}
\date{October 2, 2026}
\begin{document}
\maketitle
\begin{center}\footnotesize\bfseries
Scope notice (delivery wrapper):\ spatial geometry only;\ no PDE / Navier--Stokes / Clay regularity claim.
\end{center}
\vspace{0.75em}
\begin{abstract}
This note replaces the geometric estimate in the recovered June 19 Ring Lemma manuscript. For a three-dimensional band-limited vorticity field, a threshold relative to its $L^2$ norm gives a direction-gradient bound of order frequency to the power $5/2$. An explicit divergence-free family shows that this power cannot be reduced at a fixed threshold. A threshold relative to the $L^\infty$ norm instead gives a bound linear in frequency. These are spatial estimates; no dynamical preservation or Navier--Stokes regularity theorem is asserted. Novelty has not been audited.
\end{abstract}

Work on $\mathbb T^3=(\mathbb R/2\pi\mathbb Z)^3$ with normalized measure $dx/(2\pi)^3$. All norms below use this measure. Let $u$ be a real, mean-zero, divergence-free trigonometric polynomial, $\omega=\nabla\times u\not\equiv0$, and $\xi=\omega/|\omega|$ where $\omega\ne0$. Suppose its Fourier support lies in $|k|\le L$, $L\ge1$.

\begin{proposition}[$L^2$-threshold estimate]
For $0<c<1$, set $E_c=\{x:|\omega(x)|\ge c\|\omega\|_2\}$. Then $|E_c|>0$, and
\[
 \|\nabla\xi\|_{L^\infty(E_c)}\le \frac{C}{c}L^{5/2}.
\]
The constant depends only on the torus normalization and the frequency cutoff convention.
\end{proposition}
\begin{proof}
At nonzero vorticity,
\[
 \partial_\ell\xi=\frac{(I-\xi\otimes\xi)\partial_\ell\omega}{|\omega|},
 \qquad |\nabla\xi|\le\frac{|\nabla\omega|}{|\omega|}.
\]
Fourier Cauchy--Schwarz and Parseval give
\[
 \|\nabla\omega\|_\infty
 \le\Big(\sum_{|k|\le L}|k|^2\Big)^{1/2}\|\omega\|_2
 \le C L^{5/2}\|\omega\|_2.
\]
Dividing on $E_c$ proves the derivative bound. Also,
\[
 \|\omega\|_2^2\le c^2\|\omega\|_2^2+\|\omega\|_\infty^2|E_c|,
\]
so $|E_c|\ge(1-c^2)\|\omega\|_2^2/\|\omega\|_\infty^2>0$. This lower bound need not be uniform in $L$.
\end{proof}

\begin{proposition}[Fixed-threshold sharpness]
There are real, smooth, divergence-free velocity fields supported in a single dyadic annulus for which
\[
 \|\nabla\xi\|_{L^\infty(E_{1/2})}\ge a L^{5/2}
\]
with $a>0$ independent of $L$. Consequently a uniform $CL$ bound on the original $L^2$-threshold set is false.
\end{proposition}
\begin{proof}
Let $n\ge4$ be a power of two and define
\[
 K_n=\{(k_1,k_2,k_3):n\le k_1\le2n-1,\ 0\le k_2,k_3\le n/4\},
 \quad M=|K_n|=n(n/4+1)^2.
\]
With $P_k=I-k\otimes k/|k|^2$ and the coordinate vectors $e_2,e_3$, put
\[
 \omega_n(x)=e_3\cos(nx_1)+M^{-1/2}\sum_{k\in K_n}P_ke_2\sin(k\cdot x).
\]
This field is divergence-free and mean-zero. It is a curl: define
$\widehat u_n(k)=i k\times\widehat\omega_n(k)/|k|^2$ for $k\ne0$.
All frequencies satisfy $n\le|k|<4n$, a fixed dyadic annulus; take $L=4n$.
Orthogonality of the sine modes and the cosine gives
\[
 \|\omega_n\|_2^2=\frac12+\frac1{2M}\sum_{k\in K_n}|P_ke_2|^2\le1.
\]
At the origin, $\omega_n(0)=e_3$, so the origin lies strictly inside $E_{1/2}$. Moreover
\[
 e_2\cdot\partial_1\xi_n(0)
 =M^{-1/2}\sum_{k\in K_n}k_1\left(1-\frac{k_2^2}{|k|^2}\right)
 \ge\frac{15}{16}n\sqrt M
 \ge\frac{15}{64}n^{5/2}.
\]
Continuity on a neighborhood of the origin makes this a lower bound for the essential supremum on $E_{1/2}$. Replacing $n$ by $L/4$ gives the stated absolute constant.
\end{proof}

\begin{proposition}[Linear estimate under an additional amplitude condition]
For $0<a\le1$, define $F_a=\{x:|\omega(x)|\ge a\|\omega\|_\infty\}$. Then
\[
 \|\nabla\xi\|_{L^\infty(F_a)}\le\frac{C}{a}L.
\]
Alternatively, if $\|\omega\|_\infty\le A\|\omega\|_2$, then
\[
 \|\nabla\xi\|_{L^\infty(E_c)}\le\frac{CA}{c}L.
\]
\end{proposition}
\begin{proof}
The band-limited Bernstein estimate $\|\nabla\omega\|_\infty\le CL\|\omega\|_\infty$, combined with the direction derivative identity, gives both conclusions. It follows, for example, by convolving with a smooth Fourier cutoff equal to one on $|k|\le L$, whose differentiated kernel has $L^1$ norm at most $CL$.
\end{proof}

\begin{remark}[Scope]
The additional ratio $A$ is a spatial amplitude-control hypothesis. Frequency localization alone gives only $A\le CL^{3/2}$. Neither this hypothesis nor a bound on $F_a$ is shown here to propagate under Navier--Stokes evolution. Single-shell support is generally not invariant under nonlinear convolution. Bounds on a moving high-vorticity subset must be matched to the exact hypotheses of a geometric regularity criterion before drawing a PDE conclusion.
\end{remark}
\end{document}
